% Part: computability
% Chapter: computability-theory
% Section: fixed-point-thm

\documentclass[../../../include/open-logic-section]{subfiles}

\begin{document}

\olfileid{cmp}{thy}{fix}
\olsection{不动点定理}

再次考虑停机问题。作为临时记号，我们用 $\gn{\cfind{x}(y)}$ 表示 $\tuple{x, y}$；可将其视为值 $\cfind{x}(y)$ 的一个“名称”。利用这个记号，我们可以重述停机问题不可判定的一种证明。

问题：是否存在具有以下性质的可计算函数 $h$：对所有的 $x$ 和 $y$，
\[
h(\gn{\cfind{x}(y)}) =
\begin{cases}
1 & \text{若 $\cfind{x}(y) \fdefined$} \\
0 & \text{否则。}
\end{cases}
\]

答案：不存在；若存在，则偏函数
\[
g(x) \simeq
\begin{cases}
0 & \text{if $h(\gn{\cfind{x}(x)}) = 0$} \\
\text{无定义} & \text{否则}
\end{cases}
\]
将是可计算的，因而具有某个指标~$e$。但这样一来，就有
\[
\cfind{e}(e) \simeq
\begin{cases}
0 & \text{if $h(\gn{\cfind{e}(e)}) = 0$} \\
\text{无定义} & \text{否则，}
\end{cases}
\]
此时 $\cfind{e}(e)$ 有定义当且仅当它无定义，矛盾。

现在，观察含有 $\cfind{e}$ 的那个方程。从某种意义上说，其中存在一种自指：我们设法使 $\cfind{e}(e)$ 的值以某种方式依赖于 $\gn{\cfind{e}(e)}$。不动点定理表明，我们\emph{确实可以}一般地做到这一点——不仅仅是为了推出矛盾。

\olref{lem:fixed-equiv} 给出了表述不动点定理的两种等价方式。从逻辑上讲，既然这两个陈述都为真，它们自然是等价的；但我们真正想表达的是，它们彼此可以直接互推，因此可作为同一定理的不同表述。

\begin{lem}
\ollabel{lem:fixed-equiv}
以下陈述等价：
\begin{enumerate}
\item 对每个部分可计算函数 $g(x,y)$，存在一个指标~$e$，使得对每个~$y$，
\[
\cfind{e}(y) \simeq g(e,y).
\]
\item 对每个可计算函数~$f(x)$，存在一个指标~$e$，使得对每个~$y$，
\[
\cfind{e}(y) \simeq \cfind{f(e)}(y).
\]
\end{enumerate}
\end{lem}

\begin{proof}
$(1) \Rightarrow (2)$：给定 $f$，定义 $g$ 为 $g(x,y) \simeq
\fn{Un}(f(x),y)$。利用 (1)，可得指标~$e$，使得对每个~$y$，
\begin{align*}
\cfind{e}(y) & = \fn{Un}(f(e),y) \\
& = \cfind{f(e)}(y).
\end{align*}

$(2) \Rightarrow (1)$：给定 $g$，利用 $s$-$m$-$n$ 定理得到 $f$，使得对每个 $x$ 和~$y$，$\cfind{f(x)}(y) \simeq g(x,y)$。利用 (2)，可得指标~$e$，使得
\begin{align*}
\cfind{e}(y) & = \cfind{f(e)}(y) \\
& = g(e,y).
\end{align*}
证毕。
\end{proof}

\begin{explain}
在证明陈述 (1) 为真（从而 (2) 也为真）之前，先体会一下它的奇特之处。将 $e$ 视为一个计算机程序；陈述 (1) 是说，对于任意部分可计算的 $g(x,y)$，都能找到一个计算 $g_e(y) \simeq g(e,y)$ 的计算机程序 $e$。换言之，可以找到一个计算某个函数的程序，而该函数引用了程序自身。
\end{explain}

\begin{thm}
\olref{lem:fixed-equiv} 中的两个陈述为真。具体而言，对每个部分可计算函数 $g(x,y)$，存在一个指标~$e$，使得对每个~$y$，
\[
\cfind{e}(y) \simeq g(e,y).
\]
\end{thm}

\begin{proof}
证明所需的要素已隐含在上文关于停机问题的讨论中。令 $\fn{diag}(x)$ 为可计算函数，对每个 $x$，它返回函数 $f_x(y) \simeq \cfind{x}(x,y)$ 的一个指标，即
\[
\cfind{\fn{diag}(x)}(y) \simeq \cfind{x}(x,y).
\]
可以将 $\fn{diag}$ 视为这样一个函数：它把一个二元函数的程序变换成一个一元函数的程序，该变换将原程序固定为第一个自变元。函数 $\fn{diag}$ 可以形式化定义如下：首先，定义 $s$ 为
\[
s(x,y) \simeq \fn{Un}^2(x,x,y),
\]
其中 $\fn{Un}^2$ 是通用二元部分可计算函数的三元函数。然后，根据 $s$-$m$-$n$ 定理，存在原始递归函数 $\fn{diag}$ 满足
\[
\cfind{\fn{diag}(x)}(y) \simeq s(x,y).
\]

现在，定义函数 $l$ 为
\[
l(x,y) \simeq g(\fn{diag}(x),y).
\]
并令 $\gn{l}$ 为 $l$ 的指标。最后，令 $e = \fn{diag}(\gn{l})$。则对每个 $y$，有
\begin{align*}
\cfind{e}(y) & \simeq \cfind{\fn{diag}(\gn{l})}(y) \\
& \simeq \cfind{\gn{l}}(\gn{l}, y) \\
& \simeq l(\gn{l}, y) \\
& \simeq g(\fn{diag}(\gn{l}),y) \\
& \simeq g(e, y),
\end{align*}
符合要求。
\end{proof}

\begin{explain}
这到底是怎么回事？假设你的任务是编写一个打印自身的计算机程序。
进一步假设你使用的编程语言拥有一个丰富而奇特的字符串函数库。
具体而言，假设你的语言有一个函数 $\fn{diag}$，其工作方式如下：
给定输入字符串~$s$，$\fn{diag}$ 会找到~$s$ 中出现的每个符号 `x`，并将其替换为原字符串加引号的版本。
例如，给定字符串
\begin{quote}
\begin{verbatim}
hello x world
\end{verbatim}
\end{quote}
作为输入时，函数返回
\begin{quote}
\begin{verbatim}
hello 'hello x world' world
\end{verbatim}
\end{quote}
作为输出。在这种情况下，编写所需的程序就很容易了；可以验证
\begin{quote}
\begin{verbatim}
print(diag('print(diag(x))'))
\end{verbatim}
\end{quote}
即可实现。对于像 C++ 和 Java 这样更常见的编程语言，同样的思想（只是实现上更复杂）依然有效。

我们离不动点定理的证明只有几步之遥了。假设打印函数 $\fn{print}(x,y)$ 的一个变体接受一个字符串 $x$ 和另一个数值自变元 $y$，并将字符串 $x$ 重复打印 $y$ 次。那么“程序”
\begin{quote}
\begin{verbatim}
getinput(y);
print(diag('getinput(y); print(diag(x), y)'), y)
\end{verbatim}
\end{quote}
在输入为 $y$ 时，会将自身打印 $y$ 次。用任意函数 $g(x,y)$ 替换 $\fn{getinput}$---$\fn{print}$---$\fn{diag}$ 这一骨架，便得到
\begin{quote}
\begin{verbatim}
g(diag('g(diag(x), y)'), y)
\end{verbatim}
\end{quote}
这个程序在输入为 $y$ 时，会用程序本身和 $y$ 作为自变元来运行 $g$。将这里的“加引号”理解为“使用……的指标”，就得到了上面的证明。

现在，如果你想把证明看作形式技巧或黑魔法，这也无妨。但你应该能够复现上面给出的论证细节。当我们证明不完全性定理（以及相关的“不动点定理”）时，我们会讨论理解它为何成立的其他角度。
\end{explain}

\begin{tagblock}{lambda}
\begin{digress}
同样的思想可以用来得到一个不动点组合子。假设你有一个 λ 项 $g$，想得到另一个项 $k$，使得 $k$ β-等价于 $gk$。根据我们的记法约定定义项
\[
\fn{diag}(x) = xx
\]
和
\[
l(x) = g(\fn{diag}(x))
\]
换句话说，$l$ 就是项 $\lambd[x][g(xx)]$。令 $k$ 为项 $ll$。那么我们有
\begin{align*}
k & = (\lambd[x][g(xx)])(\lambd[x][g(xx)]) \\
& \red  g((\lambd[x][g(xx)])(\lambd[x][g(xx)])) \\
& = gk.
\end{align*}
若取
\[
Y = \lambd[g][((\lambd[x][g(xx)])(\lambd[x][g(xx)]))]
\]
则 $Yg$ 和 $g(Yg)$ 会归约到相同的项；因此 $Yg \equiv_\beta g(Yg)$。这被称为 Curry 组合子。若改为取
\[
Y = (\lambd[xg][g(xxg)])(\lambd[xg][g(xxg)])
\]
则事实上 $Yg$ 会归约到 $g(Yg)$，这是一个更强的结论。后一版本的 $Y$ 被称为 Turing 组合子。
\end{digress}
\end{tagblock}

\end{document}